% =====================================================================
%  Thesis proposal, Appendix: Skeleton residual and bubble budget
%  Owns: apx:residual
% =====================================================================

\section{Skeleton residual and bubble budget}
\label{apx:residual}

This appendix supports \cref{sec:bound}. It gives the two constants behind
\cref{eq:target} that presently exist, the lag behind the glide and the
count of passes, together with the equation that connects glide, phugoid
and skip; it says what has been computed against them, and states the range
in which each holds. The
identities are checked by computer algebra, and the leading-order
statements as limits of exact expressions, against the same definition of
the field that is evaluated numerically; each comparison with computation
was run against a failure criterion stated beforehand.

%--------------------------------------------------------------------------
\subsection{The planar glide}\label{sec:residual-glide}
%--------------------------------------------------------------------------

For spherical gravity and no rotation the altitude, speed and flight-path
angle decouple from the rest of~\crefrange{eq:rdot}{eq:psidot}:
\[
  \dot r = V\sin\gamma,\qquad
  \dot V = -\frac{D}{m} - g\sin\gamma,\qquad
  V\dot\gamma = \frac{L\cos\sigma}{m} - \Bigl(g - \frac{V^2}{r}\Bigr)\cos\gamma ,
\]
with $D/m = \rho V^2/(2\beta)$ for a ballistic coefficient~$\beta$, an
exponential density of scale height~$H_s$, and $g = \mu/r^2$. Two
abbreviations recur: the vertical lift-to-drag ratio and the net
gravitational acceleration,
\[
  E_v = \frac{L}{D}\cos\sigma, \qquad G = g - \frac{V^2}{r} .
\]
On the equilibrium glide $\gamma = 0$ and $L\cos\sigma/m = G$, so
$D/m = G/E_v$. As the speed bleeds off the glide altitude descends, and to
leading order in $\eps = H_s/R_e$ its rate is
\begin{equation}\label{eq:glide-descent}
  \dot h_{\mathrm{g}} \;=\; -\,\frac{2H_s\,g}{V E_v} .
\end{equation}

%--------------------------------------------------------------------------
\subsection{The phugoid}\label{sec:residual-phugoid}
%--------------------------------------------------------------------------

Linearized about a point of the glide with the coefficients frozen, the
three equations have the characteristic polynomial
$\lambda^3 + a_2\lambda^2 + a_1\lambda + a_0$, in which
$a_2 = 2G/(E_vV)$ exactly, $H_s\,a_1\to G$ and $a_0/a_1\to -2V/(rE_v)$ as
$H_s\to 0$. Its roots are one real root $\lambda_{\mathrm{s}}$ and a
complex pair $-\varsigma\pm i\omega$, the \emph{phugoid}: the oscillation of
altitude against flight-path angle about the glide. The sum of the roots is
$-a_2$, so $\varsigma = (a_2 + \lambda_{\mathrm{s}})/2$ exactly, and to
leading order $\lambda_{\mathrm{s}} = -a_0/a_1 = 2V/(rE_v)$. Hence
\begin{equation}\label{eq:phugoid}
  \omega^2 \;=\; \frac{G}{H_s}, \qquad
  \varsigma \;=\; \frac{g}{V E_v},
\end{equation}
each up to a relative correction of order~$\eps$. The real root is positive
and is not an instability. It is the rate at which a frozen linearization
registers that the glide itself is descending.

The two quantities in \cref{eq:phugoid} scale differently, and everything
in this appendix follows from that. The frequency grows like
$\eps^{-1/2}$. The damping rate does not contain $H_s$ at all; by
\cref{eq:glide-descent} it is $\abs{\dot h_{\mathrm{g}}}/(2H_s)$, half the
rate at which the glide descends through scale heights. The damping ratio is
therefore
\begin{equation}\label{eq:zeta}
  \zeta \;=\; \frac{\varsigma}{\omega}
  \;=\; \frac{\sqrt{g R_e}}{V E_v\sqrt{1 - V^2/(gr)}}\;\sqrt{\eps} ,
\end{equation}
of order $\sqrt{\eps}$. For $L/D = 2$ at $6$\,km/s and the terrestrial
scale height the period is about $265$\,s, the decay time
$1/\varsigma\approx 1250$\,s and $\zeta\approx 0.034$: the oscillation
lasts as long as the entry.

%--------------------------------------------------------------------------
\subsection{The residual constant}\label{sec:residual-constant}
%--------------------------------------------------------------------------

A trajectory started level on the glide altitude is not on the glide: the
glide descends at the rate \cref{eq:glide-descent} and the trajectory
trails it by one phugoid response time, ringing with an altitude amplitude
$\abs{\dot h_{\mathrm{g}}}/\omega$. In scale heights that amplitude is
$C_{\calR}\sqrt{\eps}$ with
\begin{equation}\label{eq:CR}
  C_{\calR} \;=\; \frac{2\sqrt{g R_e}}{V E_v\sqrt{1 - V^2/(gr)}} ,
  \qquad\text{so that}\qquad
  C_{\calR}\sqrt{\eps} \;=\; 2\,\zeta .
\end{equation}
The residual constant and the damping ratio are one number, read two ways.

The form has been measured rather than assumed. With the scale height
varied from $2$ to $30$\,km at fixed anchor density, the largest altitude
deviation in scale heights scales as $\eps^{0.49}$ and the largest
flight-path angle as $\eps^{0.98}$, so the altitude channel sets the bound.
Against direct integration \cref{eq:CR} is an upper bound in every case
tried, tight to about four percent. The cases vary one quantity at a time
about a generic lifting body at an angle of attack $\alpha = 16^\circ$,
$V = 6$\,km/s and zero bank, the angle of attack setting $L/D$:
\begin{center}
\begin{tabular}{@{}lcccc@{}}
\toprule
case & $L/D$ & computed & \cref{eq:CR} & ratio \\
\midrule
$\alpha = 10^\circ$ & $4.89$ & $0.813$ & $0.822$ & $0.988$ \\
$\alpha = 16^\circ$ & $3.30$ & $1.193$ & $1.215$ & $0.982$ \\
$\alpha = 25^\circ$ & $2.09$ & $1.858$ & $1.912$ & $0.972$ \\
$\alpha = 35^\circ$ & $1.41$ & $2.714$ & $2.834$ & $0.958$ \\
$V = 5.0$\,km/s & $3.30$ & $1.207$ & $1.236$ & $0.977$ \\
$V = 6.8$\,km/s & $3.30$ & $1.342$ & $1.362$ & $0.985$ \\
$\sigma = 30^\circ$ & $3.30$ & $1.373$ & $1.402$ & $0.979$ \\
$\sigma = 60^\circ$ & $3.30$ & $2.335$ & $2.422$ & $0.964$ \\
\bottomrule
\end{tabular}
\end{center}
The formula bounds rather than equals because the trajectory rings: it
tracks the envelope of the oscillation, and the computed supremum touches it
from below.

%--------------------------------------------------------------------------
\subsection{Why this is not a settling rate}\label{sec:residual-settling}
%--------------------------------------------------------------------------

\Cref{eq:CR} bounds the distance to the glide for a trajectory that starts
on it. It does not say that other trajectories approach it, and by
\cref{eq:zeta} they do so only slowly. The number of oscillations per
$e$-folding of the amplitude is $1/(2\pi\zeta)$, which grows like
$\eps^{-1/2}$. In the limit the motion normal to the glide is an undamped
oscillation on its own timescale. The glide balance is then not a normally
hyperbolic slow manifold, the hypothesis under which a slow manifold
persists and attracts~\cite{fenichel1979geometric,jones1995geometric}, and
the regular-perturbation analysis of the glide notes the same thing from
the other side: the exact solution does not approach the zeroth-order
one~\cite{lu2006asymptotic}.

The frozen rate is, moreover, not what either amplitude does. The
coefficients drift as the vehicle slows, and with that drift retained the
two amplitudes of the oscillation decay at different rates,
\begin{equation}\label{eq:decay}
  \frac{\dot A_h}{A_h} = -\frac32\,\varsigma ,
  \qquad
  \frac{\dot A_\gamma}{A_\gamma} = -\frac12\,\varsigma ,
\end{equation}
for the altitude deviation and the flight-path angle respectively: the
frequency rises as the speed falls, so the altitude swing closes faster and
the angle swing slower. Their product, the action of the oscillator, decays
at $2\varsigma$. Integration of the full equations confirms both rates to
about one percent at the terrestrial scale height.

What excites the oscillation is at least as important as what damps it. A
change of bank from $\sigma_1$ to $\sigma_2$ moves the glide altitude by
$H_s\ln(\cos\sigma_2/\cos\sigma_1)$, a fixed fraction of a scale height
however thin the layer. Under bank modulation the altitude therefore
departs from the instantaneous glide by an amount of order one in scale
heights, not of order $\sqrt{\eps}$, and so does a trajectory leaving a
pass.

Two ways of closing the bound are consistent with this.
\begin{enumerate}[leftmargin=*,itemsep=2pt]
  \item \emph{Carry the amplitude as an error.} The radius of the residual
    ball is the lag $C_{\calR}\sqrt{\eps}$ plus the ringing amplitude,
    bounded at its excitation and decaying by \cref{eq:decay}. The bound
    stays explicit and the limit system stays the glide alone, at the cost
    of a residual that is of order one in scale heights after every pass
    and every bank change.
  \item \emph{Carry the amplitude as a coordinate.} The amplitude is
    promoted to a slow coordinate of the limit system and the phase is
    averaged out. The slow variables then follow the averaged inclusion and
    the fast ones converge in the statistical sense
    of~\cite{artstein1999invariant}. The residual is the averaging error,
    which for a smooth field with a single fast phase is of order
    $\sqrt{\eps}$ on the slow timescale; a rate that is uniform over the
    selections, with its constant, is what would have to be shown. The
    limit system gains a dimension.
\end{enumerate}
The body takes the second, for a reason given in
\cref{sec:residual-oscillator}: the bank history acts on the amplitude
throughout, so a bound of the first kind does not hold uniformly over bank
histories. The first remains the form the bound takes for bank histories
that vary slowly against the oscillation.

In the unstretched state variables the two descriptions agree on the
order: an oscillation of fixed amplitude in scale heights has an altitude
excursion of order $\eps R_e$ and a flight-path-angle excursion of order
$\sqrt{\eps}$, so the occupation measures do concentrate on the glide as
$\eps\to 0$, at the rate $\sqrt{\eps}$ set by the angle. What they do not
do is concentrate at that rate in the metric of \cref{eq:target}, in which
altitude is measured in scale heights.

%--------------------------------------------------------------------------
\subsection{Glide, phugoid and skip as one oscillator}
\label{sec:residual-oscillator}
%--------------------------------------------------------------------------

Let $z = (h - h_{\mathrm{g}})/H_s$ be the altitude above the glide in scale
heights. At level flight $\ddot h = V\dot\gamma = L\cos\sigma/m - G$, and
the lift at the altitude $h_{\mathrm{g}} + H_s z$ is the lift on the glide
times $e^{-z}$, that is, $G\,e^{-z}$. To leading order in $\eps$, at frozen
speed and with the damping left out, the altitude therefore obeys
\begin{equation}\label{eq:oscillator}
  \ddot z \;=\; \omega^2\bigl(e^{-z} - 1\bigr),
  \qquad \omega^2 = \frac{G}{H_s} ,
\end{equation}
which is checked as the limit $H_s\to 0$, at fixed~$z$, of the exact
vertical acceleration of the planar field. The equation is conservative,
with energy
\[
  W \;=\; \tfrac12\,\dot z^2 + \omega^2\bigl(e^{-z} + z\bigr) ,
\]
and its potential has a single minimum, at $z = 0$. Three regimes of entry
flight are three ranges of its amplitude. At the minimum the vehicle is on
the glide. Near it the potential is quadratic, with curvature $\omega^2$,
and the motion is the phugoid of \cref{eq:phugoid}. Far from it the two
sides of the well differ. Below the glide the wall is exponential, and the
trajectory rebounds from it within a few scale heights; above, the
acceleration tends to $-\omega^2$, that is, $\ddot h\to -G$, a ballistic
arc under the net gravity. A large oscillation is therefore a skip. Let
$\gamma_0$ be the flight-path angle at which the trajectory crosses the
glide altitude, so that $\dot z = V\gamma_0/H_s$ there. The rebound lasts a
time of order $H_s/(V\gamma_0)$, the arc lasts $2V\gamma_0/G$, and the two
turning points are the roots of
\begin{equation}\label{eq:amplitude}
  e^{-z} + z - 1 \;=\; \frac{V^2\gamma_0^2}{2\,G\,H_s}
  \;=\; \frac{V^2}{2\,G\,R_e}\;\frac{\gamma_0^2}{\eps} .
\end{equation}
The amplitude in scale heights is a function of $\gamma_0/\sqrt{\eps}$,
with a factor fixed by the speed. Held fixed as $\eps\to 0$, that ratio
gives an oscillation a fixed number of scale heights deep whose period is
of order $\sqrt{\eps}$: the bounded-load limit of \cref{rem:regimes}. Sent
to infinity, it gives the skip. The classical theory describes skip paths
in these terms, as oscillatory paths of large amplitude~\cite{shi1971skip};
\cref{eq:oscillator} is the equation of which that is a statement.

What the equation leaves out is of two kinds. The damping of
\cref{eq:phugoid} and the drift of the coefficients behind \cref{eq:decay}
are each of relative order $\sqrt{\eps}$ over one period. The motion of the
equilibrium when the bank changes is not small at all. With the vertical
lift scaled by $c(t) = \cos\sigma(t)/\cos\sigma_{\mathrm{ref}}$ against a
reference bank, the equation reads $\ddot z = \omega^2(c\,e^{-z} - 1)$, its
equilibrium is at $z = \ln c$, and
\begin{equation}\label{eq:pumping}
  \dot{W} \;=\; \omega^2\,(c - 1)\,e^{-z}\,\dot z .
\end{equation}
A bank history that adds lift while the vehicle climbs and removes it while
it descends raises the energy of the oscillation at every instant, and the
opposite one lowers it. The second is what a guidance law does when it
feeds the altitude rate back into the bank command, which is how the large
phugoid oscillations of vehicles of high lift-to-drag ratio are removed in
practice~\cite{lu2014entry}. The first is open to any admissible control
that can follow the oscillation, and a reachable set has to contain what it
produces. The amplitude is a controlled variable, and that is why it is
carried as a coordinate.

%--------------------------------------------------------------------------
\subsection{Range of validity}\label{sec:residual-range}
%--------------------------------------------------------------------------

$C_{\calR}$ diverges as $V^2\to gr$. Near circular speed the lift required
vanishes, the glide altitude climbs, and by \cref{eq:phugoid} the phugoid
slows until it is no longer fast: there is no separation of scales to
exploit. The constant is therefore stated on a corridor bounded away from
circular speed, and what lies above that bound is the regime of arcs and
passes, to which the glide description does not apply. The same boundary
appears independently in the attitude dynamics, where the dynamic pressure
on the glide falls with the lift required and the trim loses stiffness, so
that a floor on dynamic pressure belongs among the hypotheses of the
corridor rather than among its consequences.

%--------------------------------------------------------------------------
\subsection{The bubble budget}\label{sec:residual-budget}
%--------------------------------------------------------------------------

By \cref{lem:energy} the specific energy is non-increasing along every
trajectory of every selection, so a pass that dissipates at least $\delta$
can occur at most $K_\delta\le\Delta\mathcal{E}/\delta$ times: this is
\cref{eq:kmax}. The dissipation of a pass is tied to how far it turns the
flight path. Where the aerodynamic force dominates gravity,
$\dd V/V = -\dd\gamma/E_v$, and a pass that turns the flight path through
$\Delta\gamma$ at constant $E_v$ leaves with
\begin{equation}\label{eq:turning}
  \frac{V_{\mathrm{out}}}{V_{\mathrm{in}}}
  \;=\; e^{-\Delta\gamma/E_v} ,
\end{equation}
the classical relation for a skip, in which
$\Delta\gamma = 2\gamma_{\mathrm{entry}}$~\cite{eggers1958comparative}. A
pass entered at $4^\circ$ with $E_v = 2$ therefore dissipates about
thirteen percent of the kinetic energy, and an entry that sheds its energy
in such passes makes about fourteen of them before the speed has fallen by
a factor $e$. The count needs no heating correlation; the
integrated heat load gives the same budget its physical reading, and the
corridor constraint on it is carried separately\apxref{rem:sutton-graves-circularity}.

How deep a pass goes is fixed in the same way. Along the planar field the
quantity $\cos\gamma - E_v H_s\,\rho/(2\beta)$ changes only through
gravity, at the rate $G\sin\gamma\cos\gamma/V$: the lift and the density
gradient cancel exactly. Where the aerodynamic force dominates it is
conserved, so a pass entered from thin air at a flight-path angle
$-\gamma_e$ levels out at the density $\rho_{\mathrm{b}}$ where
\begin{equation}\label{eq:pass-depth}
  \frac{E_v H_s\,\rho_{\mathrm{b}}}{2\beta} \;=\; 1 - \cos\gamma_e ,
  \qquad\text{that is, where}\qquad
  \frac{L\cos\sigma}{m}
  \;=\; \frac{V^2\,(1-\cos\gamma_e)}{H_s}
  \;\approx\; \frac{V^2\gamma_e^2}{2H_s} .
\end{equation}
Against gravity the load at the bottom of a pass is of order
$\gamma_e^2/\eps$, and holding it fixed as the layer thins forces
$\gamma_e\sim\sqrt{\eps}$. The same identity measures what this inner
description neglects. Across a pass the conserved quantity drifts by at
most $G\gamma_e/V$ times the time spent in the layer, which is a few times
$H_s/(V\gamma_e)$, so the correction to \cref{eq:pass-depth} is of relative
order $\eps/\gamma_e^2$: the reciprocal of the same parameter, and the
small parameter of the matching in \cref{obj:composition}.

The width of a pass has been measured as well, and against what surrounds
it. A computed pass entered at $80$\,km, $6.5$\,km/s and $-4^\circ$, which
is twice $\sqrt{\eps} = 0.034$\,rad, returns to its entry altitude after
$145$\,s, having shed $14$ percent of its kinetic energy. The arc that
follows lasts $185$\,s, against the $2V\gamma_0/G = 179$\,s of
\cref{sec:residual-oscillator}, and the period $2\pi/\omega$ of the small
oscillation about the glide at the speed of the pass is $286$\,s. These are
$0.18$, $0.23$ and $0.35$ of the slow time $\sqrt{R_e/g}$. At the
terrestrial scale height the three are not separated: the pass lasts half a
period of the oscillation, and is one of its troughs more than it is an
event short against it. An estimate that substitutes
$\sqrt{\eps}\,\sqrt{R_e/g} = 27$\,s for the width of a pass understates it
fivefold.
